\documentclass[10pt,a4paper]{amsart}
\usepackage[margin=19mm]{geometry}
\usepackage{amsmath,amssymb,mathtools}
\usepackage[scr=rsfs,cal=euler]{mathalfa}
\usepackage{xfrac}
\usepackage[unicode,hidelinks,bookmarks=false]{hyperref}
\newcommand{\ie}{i.e.\ }
\newcommand{\cf}{cf.\ }
\newcommand{\resp}{respectively\ }
\newcommand{\Subsection}{\subsection}
\providecommand{\nobreakdash}{\nobreak\hbox{-}}
\setlength{\emergencystretch}{2em}
\multlinegap0pt
\usepackage{amssymb}
\usepackage{caption}
\usepackage{float}
\usepackage{bm}
\usepackage{tikz}
\newtheorem{lemma}{Lemma}
\newcommand{\R}{\mathbb{R}}
\newcommand{\eps}{\varepsilon}
\renewcommand{\ge}{\geqslant}
\renewcommand{\le}{\leqslant}
\title{Stability of reverse isoperimetric inequalities in the plane: area, Cheeger, and inradius}
\author{Kostiantyn Drach, Kateryna Tatarko}
\date{Source: arXiv:2410.06096v2 (CC BY 4.0)}
\begin{document}
\maketitle

\noindent\emph{Source and licence.} Stability of reverse isoperimetric inequalities in the plane: area, Cheeger, and inradius. Version arXiv:2410.06096v2 (CC BY 4.0), distributed by arXiv under the Creative Commons Attribution 4.0 International licence, http://creativecommons.org/licenses/by/4.0/, as recorded in the arXiv licence metadata for this exact version and shown on the abstract page https://arxiv.org/abs/2410.06096v2. Full licence notice: \texttt{licenses/arxiv-2410.06096-CC-BY-4.0.txt}.\par\smallskip

\noindent\textit{Editorial scope.} Counted unit: the article's lemma Lem:Proximity (source0.tex lines 676--688). The article's complete original proof of it is delivered with it (source0.tex lines 690--704, one \texttt{proof} environment of 15 lines). No mathematics is written, repaired or re-derived: the statement and the proof are the article's own lines, in the article's own notation, and no retained line is substituted or re-flowed.  Retained uncounted as the source-local context the proof needs: lines 621--624.  Nothing was omitted from any retained span, so the package carries no omission record.  No retained line contains a citation, so the wrapper carries no bibliography and no bibliography entry.  No retained line contains a figure environment, an image-inclusion command or drawing code, so no image is delivered and the package declares no asset dependency.  Editorial formatting only, in addition: an AMS \texttt{amsart} preamble that restates the article's own declarations and macros used by the retained lines, namely the environments \texttt{lemma} on separate counters, together with the standard packages those lines need.  The retained environments print in this document as Lemma 1 in order of appearance, whereas the article prints the same items with the numbering of its own full document.  The complete native source of the article as distributed in the arXiv e-print for this exact version is bundled unchanged as \texttt{sources/166502/source0.tex}.
\begin{equation}
\label{Eq:LensBasic}
|L| = \frac{|\partial L|}{2} - \sin \frac{|\partial L|}{2}, \quad r(L) = 1 - \cos\frac{|\partial L|}{4}, 
\end{equation}

\begin{lemma}[Proximity of lenses]
\label{Lem:Proximity}
There exist $\eps_0 \ge 0$ and $T \ge 1$ so that the following holds for each positive $\eps \le \eps_0$: 

If $L, L' \subset \R^2$ are two $1$-convex lenses such that 
\[
1 \le \frac{r(L)}{r(L')} \le 1 + \eps,
\]
then 
\[
1 \le \frac{|\partial L|}{|\partial L'|} \le 1 + T \eps.
\]
\end{lemma}

\begin{proof}
For simplicity of notation, let us put $x := r(L)$, $x' := r(L')$. By \eqref{Eq:LensBasic}, $|\partial L|  = 4 \arccos(1 - r(L)) = 4 \arccos(1 - x) =: f(x)$. If $x' \ge \delta$ for some $\delta > 0$, then $x \ge x' \ge \delta$ and $f(x') \ge f(\delta) = 4 \arccos(1 - \delta)$. Also, $f(x) - f(x') \le Q \cdot (x - x') \le Q \eps$, where $Q$ is the Lipschitz constant of $f(x)$ for $x \in [\delta, 1]$. Therefore, we will get
\[
\frac{f(x)}{f(x')} - 1 \le \frac{Q \eps}{f(\delta)} = T\eps, \quad \text{where}\ \  T= \frac{Q}{f(\delta)}.
\]
In order to prove the result, we need to study the asymptotics of $f(x)/f(x')$ as $x' \to 0$ (and hence, as $x \to 0$). For sufficiently small $x'$, we have:
\[
\frac{f(x)}{f(x')} - 1 = \frac{f(x) - f(x')}{f(x')} \le \frac{f'(x')(x-x')}{f(x')} \le \eps \cdot \frac{f'(x')\cdot x'}{f(x')}.
\]
It is straightforward to check that $\lim_{x' \to 0}{f'(x')\cdot x'}/{f(x')} = 1/2$. Thus, there exists $\delta > 0$ so that, say, ${f'(x')\cdot x'}/{f(x')} \le 1$ for all $0 < x' < \delta$. With such $\delta$, altogether we obtain
\[
\frac{f(x)}{f(x')} - 1 \le \eps \cdot \max \left\{1, \frac{Q}{f(\delta)}\right\},
\]
so the claim follows.
\end{proof}

\end{document}
